\documentclass[12pt]{article}
\usepackage[utf8]{inputenc}
\usepackage[french]{babel}
\usepackage{amsmath}
\usepackage{enumitem}

\begin{document}

\begin{center}
    \Large \textbf{Exercices de Dénombrement}
\end{center}

\vspace{0.5cm}

\begin{enumerate}[label=\textbf{Exercice \arabic*.}]

\item On dispose des chiffres de 0 à 9.\\
Combien de nombres de 4 chiffres peut-on former ?

\vspace{0.3cm}

\item Combien de nombres de 4 chiffres \textbf{sans répétition} peut-on former ?

\vspace{0.3cm}

\item Combien de nombres de 5 chiffres contiennent \textbf{exactement 2 chiffres pairs} ?

\vspace{0.3cm}

\item Combien de nombres de 5 chiffres contiennent \textbf{2 chiffres pairs suivis de 3 chiffres impairs} ?

\vspace{0.3cm}

\item Dans une classe de 20 élèves, on choisit 5 délégués.\\
Combien de choix possibles ?

\vspace{0.3cm}

\item Dans cette même classe, on désigne :
\begin{itemize}
    \item un président
    \item un vice-président
    \item un secrétaire
\end{itemize}
Combien de possibilités ?

\vspace{0.3cm}

\item On choisit 3 élèves parmi 20, puis on leur attribue les rôles A, B et C.\\
Combien de possibilités ?

\vspace{0.3cm}

\item On choisit 3 garçons parmi 10 et 2 filles parmi 8 pour former une équipe.\\
Combien d’équipes possibles ?

\vspace{0.3cm}

\item Même question que l’exercice précédent, mais on impose que \textbf{un garçon précis} soit toujours choisi.\\
Combien d’équipes possibles ?

\vspace{0.3cm}

\item Dans une classe de 15 garçons et 10 filles, on forme une équipe de 5 personnes contenant \textbf{au moins 2 filles}.\\
Combien d’équipes possibles ?

\end{enumerate}

\vspace{0.5cm}

\textbf{Conseil :} Avant chaque exercice, se poser la question :\\
\textit{“Est-ce que je forme un groupe ou est-ce que je remplis des rôles ?”}

\end{document}